\subsection{KT-SOULE-STAB: the general-field stabilizer cutoff}
\label{sec:kt-soule-stab}

The transcription of Soul\'e's Theorem~2 and Proposition~4
in the comparison section agrees with the supplied original
PDF, pp.~779--780. In particular the coefficients are the
ordinary integers $220,210,1120,510$. There is nevertheless
a defect in the proof of its \emph{general-field} cutoff,
which does not affect the specialization to $\mathbf Q$.

\begin{proposition}[KT-SOULE-STAB]
Soul\'e's Lemma~9, pp.~777--778, asserting that a nontrivial
order-$p$ element of $\operatorname{GL}_N(\mathcal O_F)$
satisfies $p\leq1+\max([F:\mathbf Q],N)$, is false as
stated. A valid general replacement is
\[
                         p\leq N[F:\mathbf Q]+1.
                                                               \tag{SS.1}
\]
For $F=\mathbf Q$, the printed bound $p\leq N+1$ is valid.
\end{proposition}

\begin{proof}
Set $F=\mathbf Q(\sqrt5)$ and
$t=(\sqrt5-1)/2\in\mathcal O_F$, so $t^2+t-1=0$.
The matrix
\[
              M=\begin{pmatrix}0&-1\\1&t\end{pmatrix}
                \in\operatorname{SL}_2(\mathcal O_F)
\]
has characteristic polynomial $X^2-tX+1$.
Its roots are $\zeta_5,\zeta_5^{-1}$ after one embedding
of $F$; equivalently direct reduction using $t^2=1-t$
gives $M^5=I$ and $M\ne I$. It has order five, whereas
$1+\max(2,2)=3$. The equality $M^5=I$ is also checked
symbolically in the accompanying certificate.

For the replacement, regard $F^N$ as a rational vector space
of dimension $rN$, where $r=[F:\mathbf Q]$. A nonidentity
element of prime order $p$ has a primitive $p$th root of
unity as an eigenvalue; its rational minimal polynomial
therefore contains $\Phi_p$, of degree $p-1$.
Thus $p-1\leq rN$, proving (SS.1). For $r=1$ this is
the original cutoff. This argument does not require that
$\Phi_p(M)=0$ on every eigenspace, since a fixed subspace
may also be present.
\end{proof}

The precise invalid inference in the printed proof is that
$\Phi_p$ is irreducible over $F$ whenever $F$ does not
contain all $p$th roots of unity. A proper intersection
$F\cap\mathbf Q(\zeta_p)$ already makes it reducible.
The matrix example refutes that lemma, not the numerical
conclusion of Theorem~2, which might admit a separate repair.

Replacing just this step in the published homology argument
and carrying its other inputs gives the safe cutoffs
\[
 \operatorname{card}_{rN+1}
 H_q(\operatorname{GL}_N(\mathcal O_F),\mathbf Z)_{\rm tors}
       \leq h(q,N),\qquad
 \operatorname{card}_{r(2m+1)+1}
 K_m(\mathcal O_F)_{\rm tors}\leq k(m).
                                                               \tag{SS.2}
\]
Indeed every finite cell stabilizer has all prime divisors
at most $rN+1$ by Cauchy's theorem and (SS.1); the spectral
sequence and Hurewicz exclusions are then valid away from
those primes. The cell constants are unchanged. This is
a repair of the identified step conditional on the other
published cell and homology estimates, not a fresh full
proof of Soul\'e's paper. At $F=\mathbf Q$, (SS.2) gives
exactly the original cutoff $2m+2$, so the comparison used
for $K_{4k}(\mathbf Z)$ remains unchanged.

